% Generated draft: requires mathematical and source review.
\documentclass[11pt]{article}
\usepackage[margin=25mm]{geometry}
\usepackage{fontspec}
\setmainfont{Latin Modern Roman}
\usepackage{amsmath,amssymb,mathtools,unicode-math}
\setmathfont{Latin Modern Math}
\usepackage{enumitem,needspace,xurl,hyperref}
\setlength{\parindent}{0pt}
\setlength{\parskip}{5pt}
\setlength{\emergencystretch}{3em}
\newcommand{\N}{\mathbb{N}}
\newcommand{\Z}{\mathbb{Z}}
\newcommand{\Q}{\mathbb{Q}}
\newcommand{\R}{\mathbb{R}}
\newcommand{\C}{\mathbb{C}}
\newcommand{\F}{\mathbb{F}}
\newcommand{\A}{\mathbb{A}}
\newcommand{\PP}{\mathbb{P}}
\newcommand{\E}{\mathbb{E}}
\newcommand{\eps}{\varepsilon}
\newcommand{\dd}{\,\mathrm{d}}
\begin{document}
\section*{UnsolvedMath ID: 700992}
\textbf{Source problem code:} catalogue entry
\subsection{Definitions and mathematical statement}
Let \(g\in\Z\) with \(g\geq 2\). The source denotes by \(\Gamma_g\) the fundamental group of a closed surface of genus \(g\), without explicitly specifying orientability. Under the usual interpretation as a closed connected orientable surface \(\Sigma_g\), set \(\Gamma_g=\pi_1(\Sigma_g)\). The formulation below uses this interpretation; the orientability qualification requires confirmation from the source.

For every finite subset \(S\subseteq\Gamma_g\) satisfying \(\langle S\rangle=\Gamma_g\), write \(S^{-1}=\{s^{-1}:s\in S\}\). For each integer \(n\geq 0\), define
\[
 B_S(n)=\{s_1\cdots s_m:0\leq m\leq n,\ s_i\in S\cup S^{-1}\},
 \qquad \beta_S(n)=|B_S(n)|,
\]
where the empty product is the identity. Define
\[
 \omega(\Gamma_g,S)=\limsup_{n\to\infty}\beta_S(n)^{1/n},
 \qquad
 \omega(\Gamma_g)=\inf_{\substack{S\subseteq\Gamma_g\text{ finite}\\\langle S\rangle=\Gamma_g}}\omega(\Gamma_g,S).
\]
The problem is to determine the exact value of \(\omega(\Gamma_g)\) for every integer \(g\geq 2\).

As supplied background, the source states that for the free group \(F_k\) of rank \(k\geq 2\), the corresponding infimum over all finite generating sets is \(\omega(F_k)=2k-1\), citing M. Gromov, Exercise 5.13. It further states that every finite generating set of \(\Gamma_g\) contains a subset of \(2g-1\) elements generating an infinite-index subgroup whose abelianization is \(\Z^{2g-1}\), hence a free subgroup of rank \(2g-1\). The source consequently records the lower bound
\[
 \omega(\Gamma_g)\geq 4g-3.
\]
These are supplied background assertions, not an independently verified assessment or a determination of the requested infimum.
\subsection{Short English statement}
For each genus \(g\geq 2\), determine the smallest exponential growth rate obtainable from finite generating sets of the closed surface group. The supplied lower bound is \(4g-3\). The source does not explicitly specify orientability; this draft uses the usual orientable interpretation subject to confirmation. The registry labels the problem open, without independent verification here.
\subsection{Sources}
\begin{itemize}
\item Supplied local UnsolvedMath registry record: permanent ID \(700992\), source code \texttt{catalogue entry}, title \emph{Kourovka Notebook Problem 14.7}. Upstream ID: \(700992\); upstream revision: \texttt{d371ee917ac9647fbb06cd6fc529deec6f3dda82}. The record was updated on \texttt{2026-10-16}; this is a registry timestamp, not an independent literature-verification date. Its supplied status is \texttt{open}.
\item \emph{The Kourovka Notebook (Unsolved Problems in Group Theory)}, Problem 14.7, proposed by P. de la Harpe in 1999; the supplied attribution identifies the 14th issue (1999) and the 21st issue edition (2026), \texttt{arXiv:1401.0300v48}. The registry describes the collection as Issues 1--21 (1965--2026), edited by E. I. Khukhro and V. D. Mazurov.
\item M. Gromov, \emph{Structures m\'etriques pour les vari\'et\'es riemanniennes}, Cedic/F. Nathan, Paris, 1981, Exercise 5.13; supplied citation for the free-group comparison \(\omega(F_k)=2k-1\), \(k\geq 2\).
\end{itemize}
\end{document}
